\chapter{Multi-Layer Perceptron e backpropagation}
\label{cap:mlp}

\begin{inbreve}
Componendo unità non lineari in strati si ottiene il Multi-Layer Perceptron (MLP). Il capitolo lo descrive come
rete di unità e come funzione flessibile, e lo interpreta come espansione in basi adattiva. Ne discute il potere
espressivo: approssimazione universale, VC-dimension, bias induttivo. Deriva infine per intero l'algoritmo che lo
addestra, la backpropagation, che risolve il problema del \emph{credit assignment} per le unità nascoste.
\end{inbreve}

\section{Il Multi-Layer Perceptron}

Un MLP si può vedere in due modi (\figrif{fig:due-viste}):
\begin{itemize}
  \item \textbf{come rete di unità interconnesse}, organizzate in \termine{strati} (\ingl{layer}). Lo
        \textbf{strato di input} è solo la sorgente dell'input $\x$: carica il pattern, senza calcolare $\net$ né
        $f$. Lo \textbf{strato nascosto} (\ingl{hidden}) proietta sullo strato di uscita o su un altro strato
        nascosto. Lo \textbf{strato di uscita} produce $h(\x)$;
  \item \textbf{come funzione flessibile}, composta da funzioni non lineari annidate. Una rete a due strati (uno
        nascosto e uno di uscita) calcola
        \begin{equation}\label{eq:mlp}
        h(\x) = f_k\Bigg(\sum_j w_{kj}\,f_j\Big(\sum_i w_{ji}\,x_i\Big)\Bigg),
        \end{equation}
        dove $x_i$ sono le variabili di input, $w$ i parametri liberi e $f$ funzioni non lineari (sigmoidi). Si
        confronti con il modello lineare $h(\x)=\sum_iw_ix_i$: la rete è una composizione di funzioni lineari e
        non lineari.
\end{itemize}

\begin{figure}[ht]
\centering
\begin{tikzpicture}[x=1.15cm, y=0.95cm]
\foreach \i/\y in {1/2,2/1,3/0} \node[neur] (i\i) at (0,\y) {$x_\i$};
\foreach \j/\y in {1/2.5,2/1.5,3/0.5,4/-0.5} \node[neurp] (j\j) at (2.2,\y) {};
\foreach \k/\y in {1/1.6,2/0.4} \node[neurp] (k\k) at (4.4,\y) {};
\foreach \i in {1,2,3} \foreach \j in {1,2,3,4} \draw[->, grigio] (i\i) -- (j\j);
\foreach \j in {1,2,3,4} \foreach \k in {1,2} \draw[->, grigio] (j\j) -- (k\k);
\draw[->, thick] (i1) -- (j1) node[midway, above, font=\scriptsize, black] {$w_{ji}$};
\draw[->, thick] (j1) -- (k1) node[midway, above, font=\scriptsize, black] {$w_{kj}$};
\foreach \k in {1,2} \draw[->] (k\k) -- ++(0.9,0) node[right, font=\footnotesize] {$o_\k$};
\node[font=\footnotesize] at (0,-1.4) {input ($i$)};
\node[font=\footnotesize] at (2.2,-1.4) {nascosto ($j$)};
\node[font=\footnotesize] at (4.4,-1.4) {uscita ($k$)};
\node[font=\small, anchor=west] at (6.2,1.0)
  {$h(\x) = f_k\Big(\sum_j w_{kj}\,\underbrace{f_j\big(\textstyle\sum_iw_{ji}x_i\big)}_{o_j}\Big)$};
\end{tikzpicture}
\caption{Le due viste di una rete feedforward con uno strato nascosto: rete di unità (sinistra) e funzione
annidata (destra). Ogni unità nascosta e di uscita calcola $\net$ e $f$.}
\label{fig:due-viste}
\end{figure}

Una rete si descrive con il tipo di \textbf{unità} (net e funzione di attivazione), l'\textbf{architettura}
(numero di unità, topologia, numero di strati) e l'\textbf{algoritmo di apprendimento}. L'MLP standard a due
strati è \textbf{completamente connesso}; si possono avere più strati nascosti e connessioni che ``saltano''
strati.

\subsection{Notazione uniforme}

Per una generica unità $t$ (nascosta, $j$, o di uscita, $k$):
\begin{equation}\label{eq:notazione-t}
\net_t(\x) = \sum_u w_{tu}\,o_u, \qquad o_t(\x) = f_t\big(\net_t(\x)\big),
\end{equation}
dove $u$ indica una generica sorgente di input (un'unità di input $i$ o nascosta $j$). Il pattern viene caricato
nello strato di input. Si usa quindi il simbolo $o$ sia per gli input sia per le uscite delle unità: l'input
all'unità $t$ attraverso la connessione $w_{tu}$ è $o_u$. Quindi:
\begin{itemize}
  \item per i pesi $w_{ji}$ (input $\to$ nascosta) l'input è $o_i = x_i$, la componente del pattern;
  \item per i pesi $w_{kj}$ (nascosta $\to$ uscita) l'input è $o_j$, l'uscita dell'unità nascosta $j$;
  \item $w_{t0}$ è il \textbf{bias} di ogni unità $t$, con input costante $o_0=1$;
  \item unità nascosta: $o_j = f_j\big(\sum_iw_{ji}x_i\big)$; unità di uscita: $o_k = f_k\big(\sum_jw_{kj}o_j\big)$.
\end{itemize}

\subsection{Elaborazione feedforward}

L'elaborazione di un pattern procede dall'input all'uscita:
\begin{enumerate}
  \item si carica il pattern $\x$ nello strato di input;
  \item si calcolano le uscite delle unità del primo strato nascosto, poi del secondo, e così via;
  \item si calcolano le uscite dello strato di uscita, cioè $h(\x)$;
  \item si può calcolare l'errore (il delta) in uscita.
\end{enumerate}
Le reti \termine{feedforward} hanno un'unica direzione input $\to$ output. Le \textbf{reti ricorrenti} aggiungono
connessioni di feedback e quindi una memoria delle computazioni passate (capitolo~\ref{cap:rnn}).

\section{Perché le reti neurali sono flessibili}

\subsection{Task}
Lo spazio delle ipotesi è lo \textbf{spazio continuo di tutte le funzioni} rappresentabili assegnando valori ai
pesi di un'architettura data. Con unità di uscita \textbf{sigmoidali} la rete fa \textbf{classificazione}; con
uscite \textbf{lineari} fa \textbf{regressione}. Con più unità di uscita si ottengono regressione multipla o
classificatori multi-classe (ad esempio tre uscite che valgono $0{,}2$; $0{,}7$; $0{,}1$ per tre classi).

\subsection{La rete come espansione in basi adattiva}
\label{sec:lbe-adattiva}

Nella \eqref{eq:mlp} ogni termine
\[
\phi_j(\x,\w) = f_j\Big(\sum_iw_{ji}\,x_i\Big)
\]
può essere visto come calcolato da un'unità nascosta oppure come una \textbf{funzione di base} di una LBE
(\rif{sec:lbe}). Il confronto è istruttivo:
\begin{itemize}
  \item \textbf{LBE}: $h(\x) = \sum_jw_j\,\phi_j(\x)$ (regressione) o $h(\x) = f\big(\sum_jw_j\phi_j(\x)\big)$
        (classificazione), con le $\phi$ \textbf{fissate a priori}: il modello è lineare nei parametri;
  \item \textbf{rete neurale}: $h(\x) = f_k\big(\sum_jw_{kj}\,\phi_j(\x,\w)\big)$, dove le $\phi$ \textbf{dipendono dai
        pesi} e quindi \textbf{si adattano ai dati} durante l'addestramento.
\end{itemize}

\begin{idea}[Basi adattive]
Le reti neurali sono un approccio a \textbf{funzioni di base adattive}: ogni unità nascosta calcola una nuova
\textbf{feature derivata non lineare}, appresa dai dati. La capacità rappresentativa del modello è legata alla
presenza di uno strato nascosto con attivazioni non lineari. Questo strato trasforma l'input nella
\textbf{rappresentazione interna} della rete. L'apprendimento definisce una rappresentazione adatta al task. Questa
permette di estrarre le statistiche di ordine superiore utili per approssimare la funzione target.
\end{idea}

\begin{proprieta}{Le unità nascoste devono essere non lineari}
Un MLP con unità nascoste \textbf{lineari} equivale a una rete a \textbf{un solo strato}.
\end{proprieta}
\begin{proof}
Con $f$ identità nelle unità nascoste e in uscita:
\[
h(\x) = \sum_jw_{kj}\sum_iw_{ji}x_i = \sum_i\Big(\sum_jw_{kj}w_{ji}\Big)x_i = \sum_iw'_ix_i ,\qquad w'_i = \sum_jw_{kj}w_{ji},
\]
cioè ancora un modello lineare: la composizione di funzioni lineari è lineare.
\end{proof}

Il prezzo della non linearità è che il modello è \textbf{non lineare nei parametri}: l'addestramento diventa un
problema di \textbf{ottimizzazione non lineare e non convesso}.

\subsection{Approssimazione universale}
\label{sec:approssimazione-universale}

\begin{teorema}[Approssimazione universale]
Una rete con \textbf{un solo strato nascosto} (con attivazioni logistiche) e uscita lineare può approssimare
arbitrariamente bene \textbf{ogni funzione continua} su un ipercubo, purché abbia abbastanza unità nascoste. Cioè:
data $f$ e dato $\eps>0$ esiste $h$ tale che $\abs{f(\x)-h(\x)}<\eps$ per ogni $\x$ dell'ipercubo. Più in generale un
MLP può approssimare arbitrariamente bene ogni mappatura input-output, con abbastanza unità nascoste.
\end{teorema}

\begin{attenzione}[È un teorema di esistenza]
Il teorema dice che la rete \textbf{esiste}, ma non fornisce né l'\textbf{algoritmo} per trovarla né il
\textbf{numero di unità} necessario. Restano quindi due questioni: come apprendere (la backpropagation) e come
scegliere l'architettura (model selection, \rif{sec:unita-nascoste}).
\end{attenzione}

\subsection{Potere espressivo e VC-dimension}
\label{sec:vc-reti}

Il potere espressivo di una rete dipende dal \textbf{numero di unità} e dalla loro configurazione. La
VC-dimension dipende dal numero di parametri. In un MLP completamente connesso sono circa il numero di input per
il numero di unità nascoste, più le unità nascoste per le uscite. Ma la VC-dimension dipende anche dal
\textbf{valore} dei pesi:
\begin{itemize}
  \item pesi nulli: VC-dimension minima;
  \item pesi piccoli: si lavora nella parte lineare delle sigmoidi, VC-dimension piccola;
  \item pesi grandi: modello più complesso.
\end{itemize}
L'addestramento parte da pesi vicini a zero, e i pesi crescono durante il training. La rete è quindi uno
\textbf{spazio delle ipotesi di dimensione variabile}: la VC-dimension effettiva cresce durante l'apprendimento.
Su questo si basa l'early stopping (\rif{sec:early-stopping}).

Il teorema di approssimazione universale dice che uno strato nascosto basta. Non dice che basti un numero
\emph{piccolo} di unità. Esistono funzioni che con un solo strato richiederebbero un numero esponenziale di unità,
e con più strati molte meno (\rif{sec:no-flattening}).

\subsection{Il bias induttivo}

Per le reti addestrate con backpropagation il bias induttivo è legato alla \textbf{smoothness}: piccole variazioni
dell'input producono piccole variazioni dell'output (ad esempio derivata prima localmente limitata). È
un'assunzione molto comune e ragionevole: se la funzione target non è liscia (un generatore di numeri casuali) la
generalizzazione è comunque impossibile.

\section{Verso l'algoritmo di apprendimento}
\label{sec:credit-assignment}

\subsection{Il problema del credit assignment}

Il Perceptron ha un algoritmo di apprendimento, ma non rappresenta tutte le funzioni booleane. Una rete di unità a
soglia le rappresenta tutte, ma serve un algoritmo per addestrarla. Il problema è il \termine{credit assignment}:
quanta ``responsabilità'' dell'errore attribuire alle unità nascoste? Per le unità di uscita si conosce la
risposta desiderata, e quindi l'errore. Per le unità nascoste \textbf{non si sa quale dovrebbe essere la loro
uscita}: il segnale d'errore non è direttamente misurabile. In più il modello è non lineare nei pesi. Il
problema è risolto dall'algoritmo di \textbf{backpropagation}.

\begin{definizione}[Loading problem]
Data una rete e un insieme di esempi, esiste un insieme di pesi per cui la rete è consistente con gli esempi?
(Si tratta di ``caricare'' i dati di training nei parametri liberi.) Il loading problem è
\textbf{NP-completo}: non è noto un algoritmo polinomiale. In pratica però le reti si addestrano in tempi
ragionevoli con la backpropagation, senza garanzia di ottimalità.
\end{definizione}

\subsection{L'idea: estendere la discesa del gradiente}

La discesa del gradiente si estende agli MLP purché \textbf{loss e funzioni di attivazione siano differenziabili}:
bisogna trovare il delta di \textbf{ogni} unità. Il problema resta lo stesso: dati $l$ esempi $(\x_p,d_p)$ e una
loss (ad esempio l'errore quadratico), trovare $\w$ che minimizzi $E(\w) = \Remp$ calcolandone il gradiente. La
backpropagation
\begin{itemize}
  \item è semplice grazie alla \textbf{forma composizionale} del modello;
  \item usa solo quantità \textbf{locali} a ogni unità, preservandone la modularità;
  \item è \textbf{efficiente}: costa $O(\#W)$ operazioni invece di $O(\#W^2)$, dove $\#W$ è il numero di pesi. Il
        motivo è la fattorizzazione dei delta: il delta di ogni unità si calcola una volta sola e si riusa per
        tutti i suoi pesi (\rif{sec:proprieta-bp}).
\end{itemize}

\section{La backpropagation}
\label{sec:backprop}

\begin{attenzione}[Come studiare la derivazione]
La derivazione va saputa \textbf{rifare}, non memorizzare: l'obiettivo è saperla ricalcolare per loss, attivazioni o
architetture diverse.
\end{attenzione}

\subsection{Il problema}

La rete di riferimento è un MLP completamente connesso con indici $i$ (input), $j$ (nascoste), $k$ (uscita;
$K$ unità di uscita). Ogni unità di input carica la
componente del pattern: $o_i = x_i$. Si scrive $TR = \{(\x_1,\dd_1),\dots,(\x_l,\dd_l)\}$ con $\dd_p$ vettore dei
target delle $K$ uscite. L'\textbf{errore totale} è
\begin{equation}\label{eq:etot}
E_{tot} = \sum_pE_p,\qquad E_p = \frac12\sum_{k=1}^{K}(d_k-o_k)^2 ,
\end{equation}
dove $E_p$ è l'errore sull'esempio $p$ su tutte le uscite; il fattore $\frac12$ serve solo a semplificare la
derivata. D'ora in poi l'indice $p$ è omesso: $o_k$ sta per $o_k(\x_p)$ e $d_k$ per $d_{p,k}$.

\begin{definizione}[Il problema della backpropagation]
Stimare il \textbf{contributo delle unità nascoste all'errore} in uscita, cioè calcolare per ogni peso della rete
$\partial E_{tot}/\partial w$: la \termine{regola delta generalizzata} (\ingl{Generalized Delta Rule}).
\end{definizione}

L'obiettivo è trovare i pesi che minimizzano $E_{tot}$ scendendo lungo la sua superficie seguendo il gradiente. A
differenza del modello lineare la superficie \textbf{non è convessa} (\figrif{fig:superficie-nonconvessa}).

\begin{figure}[ht]
\centering
\begin{tikzpicture}
\begin{axis}[width=7.6cm, height=5.4cm, view={40}{38}, axis lines=none, colormap={gr}{gray(0cm)=(0.93); gray(1cm)=(0.5)}]
\addplot3[surf, shader=faceted interp, faceted color=black!40, line width=0.1pt, domain=-3:3, y domain=-3:3, samples=31, samples y=31]
  {1.8 - 1.3*exp(-((x-1)^2+(y-1)^2)/0.8) - 0.9*exp(-((x+1.4)^2+(y+0.8)^2)/0.6) - 0.6*exp(-((x+0.8)^2+(y-1.8)^2)/0.4) + 0.08*(x^2+y^2)};
\addplot3[only marks, mark=*, mark size=1.8pt, rosso] coordinates {(1,1,0.66)};
\node[font=\footnotesize, rosso] at (axis cs:1.6,1.4,0.9) {minimo};
\node[font=\footnotesize] at (axis cs:3.4,0,1.8) {$w_1$};
\node[font=\footnotesize] at (axis cs:0,-3.5,1.8) {$w_2$};
\end{axis}
\end{tikzpicture}
\caption{Una superficie d'errore non convessa nello spazio di due pesi (andamento qualitativo): più minimi
locali, e il punto di arrivo della discesa dipende dal punto di partenza.}
\label{fig:superficie-nonconvessa}
\end{figure}

\begin{algoritmo}[Backpropagation (schema iterativo)]
Inizializzare tutti i pesi $\w$ e $\eta$; calcolare le uscite delle unità ed $E_{tot}$.\\
Finché $E_{tot}>\epsilon$ (o un altro criterio di arresto):
\begin{enumerate}
  \item per ogni peso $w$: $\Delta w = -\dfrac{\partial E_{tot}}{\partial w}$ \hfill (passo 1: calcolo del gradiente)
  \item $w_{\mathit{new}} = w + \eta\,\Delta w\ (+\,\text{eventuali termini aggiuntivi})$ \hfill (passo 2: regola di aggiornamento)
  \item ricalcolare le uscite ed $E_{tot}$.
\end{enumerate}
\end{algoritmo}

Bisogna distinguere il \textbf{calcolo del gradiente} (passo 1, ciò che fa la backpropagation) dalla
\textbf{regola di aggiornamento} (passo 2). La regola standard è $w+\eta\,\Delta w$. I termini aggiuntivi del
passo 2 sono quelli delle varianti: il momentum, la penalità del weight decay, oppure regole diverse come
Quickprop e R-prop (capitolo~\ref{cap:addestramento}).
Per il passo 1:
\[
\Delta w = -\frac{\partial E_{tot}}{\partial w} = -\sum_p\frac{\partial E_p}{\partial w} \overset{\mathrm{def}}{=} \sum_p\Delta_pw .
\]

\subsection{La derivazione}

\subsubsection{Il gradiente per un peso generico}
Si consideri un peso generico $w_{tu}$, che porta all'unità generica $t$ l'input proveniente dall'unità generica
$u$ (che può essere un $i$ o un $j$). Poiché $\net_t = \sum_sw_{ts}\,o_s$ e $o_t = f_t(\net_t)$, la regola della
catena attraverso $\net_t$ dà
\begin{equation}\label{eq:bp-generico}
\Delta_pw_{tu} = -\frac{\partial E_p}{\partial w_{tu}}
= \underbrace{-\frac{\partial E_p}{\partial\net_t}}_{\delta_t}\cdot\underbrace{\frac{\partial\net_t}{\partial w_{tu}}}_{o_u}
= \delta_t\,o_u ,
\end{equation}
perché $\frac{\partial\net_t}{\partial w_{tu}} = \frac{\partial\sum_sw_{ts}o_s}{\partial w_{tu}} = o_u$: tutti i
termini della somma sono costanti rispetto a $w_{tu}$ tranne quello con $s=u$. Si è definito il \termine{delta
dell'unità $t$}
\[
\delta_t = -\frac{\partial E_p}{\partial\net_t},
\]
che si espande ancora con la regola della catena passando per $o_t = f_t(\net_t)$:
\begin{equation}\label{eq:delta-generico}
\delta_t = -\frac{\partial E_p}{\partial\net_t} = -\frac{\partial E_p}{\partial o_t}\cdot\frac{\partial o_t}{\partial\net_t} = -\frac{\partial E_p}{\partial o_t}\cdot f'_t(\net_t).
\end{equation}
Resta da calcolare $-\partial E_p/\partial o_t$, e qui si distinguono due casi.

\subsubsection{Caso 1: unità di uscita ($t=k$)}
L'uscita $o_k$ compare direttamente in $E_p$:
\[
-\frac{\partial E_p}{\partial o_k} = -\frac{\partial\,\frac12\sum_{r=1}^{K}(d_r-o_r)^2}{\partial o_k} = (d_k-o_k),
\]
perché solo il termine $r=k$ dipende da $o_k$, con derivata $\frac12\cdot2(d_k-o_k)\cdot(-1)$. Quindi
\begin{equation}\label{eq:delta-uscita}
\boxed{\;\delta_k = (d_k-o_k)\,f'_k(\net_k)\;}
\end{equation}
È il delta della singola unità sigmoidale (\rif{sec:lms-sigmoide}): l'errore è misurabile perché si conosce il
target.

\subsubsection{Caso 2: unità nascosta ($t=j$)}
Un'unità nascosta non ha target e $o_j$ non compare direttamente in $E_p$. Però $o_j$ influenza $E_p$
\textbf{attraverso tutte le unità di uscita} a cui è collegata, perché ogni $\net_k$ dipende da $o_j$. La regola
della catena con una somma su tutte le $k$ dà
\[
-\frac{\partial E_p}{\partial o_j} = \sum_{k=1}^{K}\underbrace{-\frac{\partial E_p}{\partial\net_k}}_{\delta_k}\cdot\frac{\partial\net_k}{\partial o_j} = \sum_{k=1}^{K}\delta_k\,w_{kj},
\]
poiché $\frac{\partial\net_k}{\partial o_j} = \frac{\partial\sum_sw_{ks}o_s}{\partial o_j} = w_{kj}$. Quindi
\begin{equation}\label{eq:delta-nascosta}
\boxed{\;\delta_j = \Big(\sum_{k=1}^{K}\delta_k\,w_{kj}\Big)\,f'_j(\net_j)\;}
\end{equation}

\begin{idea}[Il punto cruciale]
Il termine $-\partial E_p/\partial\net_k$ è \textbf{già stato calcolato}: è $\delta_k$. Il delta di un'unità
nascosta si ottiene \textbf{propagando all'indietro} i delta delle unità di uscita. Sono pesati con gli stessi
pesi $w_{kj}$ usati in avanti. L'errore ``risale'' la rete, da cui il nome \emph{back-propagation}. Questo
risolve il credit assignment. La ``colpa'' di un'unità nascosta è la somma delle colpe delle unità a cui
contribuisce, pesata da quanto contribuisce (\figrif{fig:retropropagazione}).
\end{idea}

\begin{figure}[ht]
\centering
\begin{tikzpicture}[x=1.4cm, y=1.05cm]
\node[neur] (i) at (0,0.5) {$i$};
\node[neurp] (j) at (2,0.5) {$j$};
\foreach \k/\y in {1/2.0,2/0.5,3/-1.0} \node[neurp] (k\k) at (5,\y) {$k_\k$};
\draw[->] (i) -- (j) node[midway, above, font=\footnotesize] {$w_{ji}$};
\foreach \k in {1,2,3} \draw[->] (j) -- (k\k) node[pos=0.72, above, font=\scriptsize, sloped] {$w_{k_\k j}$};
\foreach \k in {1,2,3} \node[font=\scriptsize, rosso, below right] at (k\k.south west) {$\delta_{k_\k}$};
\draw[->, rosso, densely dashed, very thick] (4.6,2.75) to[bend right=30] node[pos=0.45, above left, font=\scriptsize, rosso, align=right] {retropropagazione dei $\delta_k$} (j.north);
\node[font=\footnotesize, rosso] at (2,-0.65) {$\delta_j = \big(\sum_k\delta_kw_{kj}\big)f'_j(\net_j)$};
\foreach \k in {1,2,3} \draw[->] (k\k) -- ++(0.9,0) node[right, font=\scriptsize] {$o_{k_\k}$ vs $d_{k_\k}$};
\end{tikzpicture}
\caption{Retropropagazione dei delta: l'influenza dell'unità nascosta $j$ sull'errore si calcola dai delta delle
unità a cui è collegata e dai pesi $w_{kj}$. Ogni termine coinvolge solo unità e pesi adiacenti.}
\label{fig:retropropagazione}
\end{figure}

Lo stesso risultato si può ottenere derivando direttamente la definizione di $E_p$ (per esercizio). La forma
vista però è più istruttiva, perché mostra il flusso delle derivate parziali \textbf{sulla rete}.

\subsection{Sintesi delle formule}

\begin{table}[ht]
\centering\small
\caption{La backpropagation per una rete a due strati: $\Delta_pw_{tu} = \delta_t\,o_u$.}
\begin{tabular}{lll}
\toprule
Peso & Aggiornamento & Delta \\
\midrule
$w_{kj}$ (nascosta $\to$ uscita) & $\Delta_pw_{kj} = \delta_k\,o_j$ & $\delta_k = (d_k-o_k)\,f'_k(\net_k)$\\
$w_{ji}$ (input $\to$ nascosta) & $\Delta_pw_{ji} = \delta_j\,o_i$ & $\delta_j = \big(\sum_k\delta_kw_{kj}\big)f'_j(\net_j)$\\
\bottomrule
\end{tabular}
\end{table}

$\delta_t$ è il \textbf{segnale d'errore} disponibile all'unità $t$ e $o_u$ l'input a $t$ proveniente da $u$.
La derivazione vale per \textbf{qualsiasi numero di strati nascosti}: i delta di un'unità arrivano da tutte le
unità a cui è connessa,
\begin{equation}\label{eq:delta-generale}
\delta_j = \Big(\sum_{h\in\mathrm{successori}(j)}\delta_h\,w_{hj}\Big)\,f'_j(\net_j),
\end{equation}
e la regola di aggiornamento per ogni esempio è $w_{tu}^{\mathit{new}} = w_{tu} + \eta\,\delta_t\,o_u$. Usando
$\Delta_pw$ esempio per esempio si ottiene la versione \textbf{on-line}; sommando i contributi di tutti gli esempi
prima di aggiornare, la versione \textbf{batch}.

\begin{algoritmo}[Ciclo di addestramento completo]
\begin{enumerate}
  \item \textbf{Fase forward}: calcolare le uscite di tutte le unità, strato per strato.
  \item Calcolare errori e delta nello strato di uscita.
  \item \textbf{Fase backward}: propagare i delta dallo strato di uscita agli strati nascosti.
  \item Aggiornare i pesi.
  \item Ricominciare, fino al criterio di arresto.
\end{enumerate}
\end{algoritmo}

\begin{attenzione}[Non dimenticare i bias]
Calcolo e aggiornamento valgono anche per i \textbf{bias} di tutte le unità: con input costante $o_0=1$ e bias
$w_{t0}$, si ha $\Delta w_{t0} = \delta_t\cdot1$.
\end{attenzione}

\subsection{Proprietà e interpretazione}
\label{sec:proprieta-bp}

\begin{itemize}
  \item \textbf{Calcoli solo locali}: ogni unità usa solo informazioni delle unità vicine (sotto e sopra). Il resto
        arriva indirettamente con il flusso dei delta. Eppure la modifica di un peso ha un \textbf{effetto
        globale} sulla rete. L'influenza complessa di ogni $w_{tu}$ sull'errore globale è stata
        \textbf{scomposta} in componenti semplici, i delta locali.
  \item \textbf{Processo di diffusione}: una propagazione locale dai ``genitori'' dell'unità all'unità e
        dall'unità ai ``figli''.
  \item \textbf{Efficienza (la fattorizzazione dei delta)}: il delta $\delta_t$ di ogni unità si calcola \textbf{una
        sola volta} e si riusa per \textbf{tutti} i suoi pesi $w_{tu}$. Il costo è proporzionale al numero di pesi,
        non al suo quadrato: con $10^9$ pesi sono $10^9$ operazioni invece di $10^{18}$.
\end{itemize}

\begin{idea}[Tre livelli di lettura]
\begin{enumerate}
  \item I singoli \textbf{passaggi matematici}.
  \item Il significato delle derivate parziali come \textbf{variazioni locali} di una quantità rispetto a
        un'altra.
  \item Il \textbf{quadro generale}: la scomposizione serve a trovare i delta di ogni livello, scomponendo
        l'errore \emph{sulla rete} che si ha davanti.
\end{enumerate}
\end{idea}

\begin{attenzione}[Terminologia]
Non si dice ``rete backprop'': la rete neurale è il \textbf{modello}, la backpropagation è un \textbf{algoritmo di
addestramento} basato sul calcolo del gradiente attraverso la rete.
\end{attenzione}

\begin{esempio}[Un passo di backpropagation]
Rete con un input $x=1$, un'unità nascosta sigmoidale e un'unità di uscita lineare, senza bias; pesi
$w_{ji}=0{,}5$, $w_{kj}=1$, target $d=1$, $\eta=0{,}1$.
\begin{itemize}
  \item \textbf{Forward}: $\net_j=0{,}5$, $o_j = \sigma(0{,}5)\approx0{,}622$; $o_k = 1\cdot0{,}622 = 0{,}622$.
  \item \textbf{Delta di uscita} (lineare, $f'=1$): $\delta_k = (1-0{,}622)\cdot1 = 0{,}378$.
  \item \textbf{Delta nascosto}: $\delta_j = (\delta_kw_{kj})\,\sigma'(\net_j) = 0{,}378\cdot0{,}622\cdot(1-0{,}622)\approx0{,}089$.
  \item \textbf{Aggiornamenti}: $\Delta w_{kj} = \eta\,\delta_k\,o_j\approx0{,}1\cdot0{,}378\cdot0{,}622\approx0{,}0235$;
        $\Delta w_{ji} = \eta\,\delta_j\,x\approx0{,}0089$.
\end{itemize}
Entrambi i pesi aumentano, facendo crescere l'uscita verso il target.
\end{esempio}

\begin{puntichiave}
\begin{itemize}
  \item MLP: $h(\x) = f_k\big(\sum_jw_{kj}f_j(\sum_iw_{ji}x_i)\big)$; una LBE con basi $\phi_j(\x,\w)$ adattive.
  \item Unità nascoste lineari $\Rightarrow$ modello lineare. Un solo strato nascosto basta per approssimare
        funzioni continue (teorema di esistenza).
  \item La VC-dimension dipende anche dal valore dei pesi: spazio delle ipotesi di dimensione variabile.
  \item Backpropagation: $\Delta_pw_{tu} = \delta_to_u$, $\delta_k = (d_k-o_k)f'_k$,
        $\delta_j = (\sum_k\delta_kw_{kj})f'_j$; costo lineare nel numero di pesi.
\end{itemize}
\end{puntichiave}

\begin{domande}
\begin{enumerate}
  \item Perché le unità nascoste devono essere non lineari?
  \item In che senso una rete neurale è una LBE adattiva?
  \item Enunciare il teorema di approssimazione universale e discuterne i limiti.
  \item Cos'è il problema del credit assignment? E il loading problem?
  \item Derivare la backpropagation per un MLP a due strati con errore quadratico. Come cambia con una loss diversa
        o con più strati nascosti?
  \item Perché la backpropagation ha costo lineare nel numero dei pesi?
\end{enumerate}
\end{domande}
